%% Source: MANUSCRIPT.md "Materials and Methods > Multi-source thermodynamics"
%%         (drafts/methods_multi_source_thermodynamics.md, inlined).
\subsection{Multi-source thermodynamics}\label{sec:methods-thermo}

The 2020 release combined two thermodynamics sources: eQuilibrator as the
primary estimator of \dfG{} and \drG{}, and a group-contribution (GC) method as
a fallback for compounds and reactions eQuilibrator could not estimate. For this
update we integrate four sources, each providing an independent estimate
wherever its coverage permits.

\paragraph{eQuilibrator refresh.} \TBD[version of eQuilibrator] was run against
the 2026 compound set at pH 7.0, ionic strength 0.25\,M and temperature
298.15\,K, matching the 2020 conditions. Coverage: \TBD[N compounds with
$\Delta_{\mathrm{f}}G'$, N reactions with accepted $\Delta_{\mathrm{r}}G'$].

\paragraph{Group Contribution rebuilt from MFAToolkit.} The 2020 fallback GC
method has been rebuilt from the MFAToolkit implementation against the current
compound set, so the group definitions and parameter fits reflect the enlarged
compound base rather than the 2020 snapshot. Coverage: \TBD{}.

\paragraph{dGpredictor retrained on ModelSEED.} dGpredictor (an ML-based
$\Delta G$ estimator) has been retrained using the ModelSEED-curated compound
structures as training data. This produces a ModelSEED-native estimator rather
than one trained on a competing biochemistry database, and lets us report
$\Delta G$ estimates for compounds the other three methods cannot cover.
Training procedure and hyperparameters: \TBD{}. Coverage: \TBD{}.

\paragraph{OpenTECR integration.} OpenTECR is a community effort to standardize
experimentally-measured thermodynamic constants; we integrate its values where a
mapping between OpenTECR reaction identifiers and ModelSEED reaction IDs exists.
Coverage: \TBD{}.

\paragraph{Direction assignment.} For each reaction with an accepted \drG{}
(from any source, chosen with per-compound-consistency constraints described in
the 2020 paper), we apply the same reversibility rules used in 2020 (based on
Jankowski \textit{et al.}~\cite{jankowski2008}) to assign one of \texttt{=}
(reversible), \texttt{\textgreater} (irreversible left-to-right), or
\texttt{\textless} (irreversible right-to-left). On top of this we add
per-reaction-class heuristic overlays --- hard rules for classes where the
standard thermodynamic assignment is known to disagree with textbook biology.
Rule set: \TBD[discussed in the working session but not yet codified in a repo
script]. Overlays are applied after the base assignment and their scope is
documented per class.

\paragraph{Combined direction ledger.} For every reaction, the pipeline records:
the \drG{} estimate from each of the four sources (with per-source uncertainty),
the base direction assignment from each source, the heuristic overlay if
applicable, and the final direction. This ledger is the input to the
reaction-direction sensitivity study reported in
Section~\ref{sec:results-direction}.
